\section{Related Work}
\label{sec:related}

% ═══════════════════════════════════════════════════════════════════════════════
% 2.1  AMP classification and potency prediction
% ═══════════════════════════════════════════════════════════════════════════════
\subsection{AMP Classification and Potency Prediction}

Computational AMP discovery began with supervised classifiers.
AMPlify~\citep{amplify}, AMP Scanner~\citep{ampscanner}, and
iAMPpred~\citep{iamppred} demonstrated that recurrent, convolutional, and
classical machine-learning models can discriminate antimicrobial from
non-antimicrobial sequences with high accuracy.
BERT-style peptide encoders later pushed supervised screening further by
introducing contextual sequence representations~\citep{ampbert,peptidebert}.
These classifiers answer a binary question---active or not---but
real-world peptide development also requires quantitative potency estimation.
GRAMPA~\citep{grampa} and DBAASP~\citep{dbaasp3} provide large multi-species
MIC measurements, and the recent QMAP benchmark~\citep{qmap2026} defines
precomputed homology-aware splits that control for sequence similarity between
training and test sets, addressing a long-standing comparability problem in AMP
regression.
Our work uses the same frozen encoder for both binary classification and
quantitative MIC prediction, and evaluates under both standard and
homology-aware protocols to test how far a single representation can stretch.

% ═══════════════════════════════════════════════════════════════════════════════
% 2.2  AMP generation
% ═══════════════════════════════════════════════════════════════════════════════
\subsection{AMP Generation}

Early generative approaches to AMP design used variational
autoencoders~\citep{pepvae} and generative adversarial
networks~\citep{gupta2018}, followed by more recent work combining deep
generative models with cell-free biosynthesis~\citep{pandi2023}, diffusion
frameworks~\citep{diffamp}, and knowledge-aware prompt
conditioning~\citep{kppepgen}.
HydrAMP~\citep{hydramp} showed that conditional generation can propose
peptides with improved predicted potency and reduced hemolysis.
A common thread is that generators are evaluated primarily by aggregate
statistics---average hydrophobicity, net charge, or predicted AMP score of
a generated batch---rather than by asking which specific target properties
the model can and cannot control.
Our generation experiments adopt a different evaluation stance: we condition
on individual physicochemical properties and MIC targets, then audit control
success property by property, explicitly mapping where generation succeeds
and where physical or data constraints set the boundary.
We also compare three decoding paradigms---autoregressive, non-autoregressive,
and masked diffusion---on the same frozen encoder, isolating the effect of
decoding strategy from the effect of representation quality.

% ═══════════════════════════════════════════════════════════════════════════════
% 2.3  Self-supervised representation learning for proteins and peptides
% ═══════════════════════════════════════════════════════════════════════════════
\subsection{Self-Supervised Representation Learning for Proteins and Peptides}

Large protein language models such as ESM-2~\citep{esm2} and
ProtTrans~\citep{prottrans} learn general sequence representations from
evolutionary-scale corpora and transfer well to many protein tasks.
UniRep~\citep{unirep} earlier showed that unsupervised sequence pre-training
can support downstream protein prediction even with simpler architectures.
At the peptide scale, domain-specific models---PepBERT~\citep{pepbert},
LLAMP~\citep{llamp}, PeptideBERT~\citep{peptidebert}---have demonstrated
that pre-training on curated peptide corpora can match or exceed general
protein models at a fraction of the parameter count, establishing the value
of domain-specific representation learning for short sequences.

These models span a wide range of scales and training data, but they share
a common pre-training paradigm: predicting masked or next tokens from
sequence context.
An alternative family of self-supervised objectives operates entirely in
representation space.
The Joint-Embedding Predictive Architecture (JEPA)~\citep{jepa_assran},
originally proposed for images, predicts the \emph{embedding} of a target
region from a context region, encouraging the encoder to capture abstract
structure rather than reconstruct surface-level inputs.
To our knowledge, JEPA-style representation-level pre-training has not been
applied to biological sequences.
\projectname{} tests this objective on AMPs---a domain where activity depends
on holistic physicochemical patterns rather than specific residue
identities---and compares the resulting representations against both a
general-purpose protein model (ESM-2) and an identically architected
MLM-pre-trained encoder.
